\begin{textsample}{NEG Dim 1 – global\_south\_2025\_03 – Score 1.00 – global\_south\_2025\_03/121883736.txt}  \label{ex:f1_neg_166}
The Israeli government says it has blocked the entry of all humanitarian aid into Gaza because the first phase of the ceasefire with Hamas has expired.
Prime Minister Benjamin Netanyahu ' s office said Hamas had so far refused to accept a temporary ceasefire extension under a proposal by US President Donald Trump ' s envoy, Steve Witkoff.
A Hamas spokesman called the move"cheap blackmail"and a"coup"on the ceasefire agreement and urged mediators to get Israel to resume the supply of aid.
The Palestinian group wants phase two of the deal to go ahead as originally negotiated, with the release of hostages and Palestinian prisoners and the withdrawal of Israeli forces from Gaza.
On Friday night, Hamas said it would not agree to any extension of phase one without guarantees from US, Qatari and Egyptian mediators that phase two would eventually take place.
A statement from Netanyahu ' s office said :"With the end of Phase 1 of the hostage deal, and in light of Hamas ' s refusal @ @ @ @ @ @ @ @ @ @ which Israel agreed - Prime Minister Netanyahu has decided that, as of this morning, all entry of goods and supplies into the Gaza Strip will cease."Israel will not allow a ceasefire without the release of our hostages.
If Hamas continues its refusal, there will be further consequences."The Hamas spokesman said :"Netanyahu ' s decision to stop aid going into Gaza once again shows the ugly face of the Israeli occupation...
The international community must apply pressure on the Israeli government to stop starving our people."Late last night, Netanyahu ' s office said Israel had agreed to a US proposal for the ceasefire to continue for about six weeks during the Muslim Ramadan and Jewish Passover periods.
If, at the end of this period, negotiations reached a dead end, Israel would reserve the right to go back to war.
US envoy Witkoff has not made his proposal public.
According to Israel, it would begin with the release of half of all the remaining living @ @ @ @ @ @ @ @ @ @ would immediately start negotiations if Hamas changed its position on the six-week ceasefire extension.
The first phase of the ceasefire that came into force on 19 January expired on Saturday.
It halted 15 months of fighting between Hamas and the Israeli military, allowing the release of 33 Israeli and five Thai hostages for about 1,900 Palestinian prisoners and detainees.
But negotiations on phase two, including the release of all remaining living hostages and the withdrawal of Israeli troops from Gaza, have barely begun.
There are believed to be 24 hostages alive, with another 39 presumed to be dead.
Hamas carried out an unprecedented attack on Israel on 7 October 2023, killing about 1,200 people and taking another 251 hostage.
Israel responded with an air and ground campaign in the Gaza Strip, during which at least 48,365 people have been killed, according to the territory ' s Hamas-run health ministry.
DISCLAIMER : The Views, Comments, Opinions, Contributions and Statements made by Readers and Contributors on this platform do @ @ @ @ @ @ @ @ @ @ Limited.

% matched lemmas: 
\end{textsample}
